% =====================================================================
%  ABSTRAK (Indonesia) + ABSTRACT (English)
% =====================================================================

\chapter*{Abstrak}
\addcontentsline{toc}{chapter}{Abstrak}

\noindent\textbf{[JUDUL TESIS DALAM BAHASA INDONESIA]}

\medskip\noindent\textbf{Nama:} [Nama Mahasiswa] \quad \textbf{NPM:} [XXXXXXXXXX]

\medskip
\noindent
[Isi abstrak dalam Bahasa Indonesia. Panjang 200--300 kata. Abstrak mencakup: (1) latar
belakang dan motivasi penelitian, (2) tujuan penelitian, (3) metode/pendekatan yang
digunakan, (4) hasil utama yang diperoleh, (5) kesimpulan dan kontribusi. Tidak mengandung
rujukan pustaka, singkatan tanpa penjelasan, atau persamaan matematis.]

\medskip
\noindent\textbf{Kata kunci:} [kata kunci 1]; [kata kunci 2]; [kata kunci 3]; [kata kunci 4]; [kata kunci 5]

\clearpage

% ── Abstract (English) ───────────────────────────────────────────────
\chapter*{Abstract}
\addcontentsline{toc}{chapter}{Abstract}

\noindent\textbf{[THESIS TITLE IN ENGLISH]}

\medskip\noindent\textbf{Name:} [Student Name] \quad \textbf{SID:} [XXXXXXXXXX]

\medskip
\noindent
[Abstract content in English. Same structure as the Indonesian abstract: (1) background
and motivation, (2) research objectives, (3) methods/approach, (4) main results,
(5) conclusions and contributions. 200--300 words. No references, unexplained acronyms,
or mathematical equations.]

\medskip
\noindent\textbf{Keywords:} [keyword 1]; [keyword 2]; [keyword 3]; [keyword 4]; [keyword 5]

\clearpage
